\uuid{A9cK}
\titre{Matrice nilpotente/Inversibilité}
\auteur{hanani ; mzari}
\organisation{exo7}
\licence{CC BY-NC-SA 4.0}
\datecreate{2020-01-20}
\niveau{L1}
\module{Algèbre}
\chapitre{Matrice}
\sousChapitre{Inverse, méthode de Gauss}
\difficulte{3}
\theme{Inverse d'une matrice, Matrice nilpotente, Inversibilité}
\exo7qcmid{lille-2.matrices.034}

\contenu{
\begin{question}
\qclassification{108.05}
\qtags{motcle=Matrice nilpotente/Inversibilité}
Soit $A$ une matrice nilpotente, c.à.d il existe un entier $n\ge 1$ tel que $A^n=0$. On notera $I$ la matrice identité. Quelles sont les assertions vraies ?
\begin{answers}
\bad{$A$ est inversible}.
\bad{$A$ est inversible et $A^{-1} = A^{n-1}$}.
\good{Il existe $a\in \Rr$, tel que $A-aI$ n'est pas inversible}.
\good{Pour tout $a\in \Rr^*$, $A-aI$ est inversible}.
\end{answers}
\begin{explanations} On suppose que $A$ est inversible, alors $A$ est non nul et il existe une matrice $C$ telle que $AC=I$.
Soit $m$ le plus petit entier $\ge 1$ tel que $A^m=0$. Alors, $0=A^mC=A^{m-1}(AC)=A^{m-1}$, ce qui est absurde.
Par conséquent, $A$ n'est pas inversible.\\
Comme $A$ n'est pas inversible, pour $a=0$, $A-aI$ n'est pas inversible.\\
Soit $a\in \Rr^*$, de l'égalité : $(A-aI)(a^{n-1}I+a^{n-2}A+\dots + aA^{n-2}+A^{n-1})=A^n-a^nI=-a^nI$, on déduit que
$A-aId$ est inversible et que
$\displaystyle (A-aId)^{-1}=-\frac{1}{a^n}\left(a^{n-1}I+a^{n-2}A+\dots + aA^{n-2}+A^{n-1}\right)$.
\end{explanations}
\end{question}
}
